\documentclass[UTF8,fontset=none,11pt,a4paper]{ctexart}
\usepackage[left=21mm,right=21mm,top=27mm,bottom=22mm,headheight=26pt,headsep=6mm]{geometry}
\usepackage{fontspec}
\setmainfont{texgyretermes-regular.otf}[BoldFont=texgyretermes-bold.otf,ItalicFont=texgyretermes-italic.otf]
\setCJKmainfont{Songti.ttc}[Path=/System/Library/Fonts/Supplemental/,FontIndex=6,AutoFakeBold=2]
\setCJKsansfont{Songti.ttc}[Path=/System/Library/Fonts/Supplemental/,FontIndex=6,AutoFakeBold=2]
\usepackage[table]{xcolor}
\definecolor{TableHead}{HTML}{EEF0ED}
\definecolor{Muted}{HTML}{666C68}
\usepackage{booktabs,longtable,array,fancyhdr,needspace,xurl}
\usepackage[unicode,hidelinks]{hyperref}
\urlstyle{same}
\linespread{1.32}
\setlength{\parindent}{2em}
\setlength{\parskip}{3pt}
\setlength{\emergencystretch}{3em}
\setlength{\LTpre}{8pt}\setlength{\LTpost}{10pt}
\ctexset{section={format=\fontsize{15}{20}\selectfont\bfseries,beforeskip=18pt,afterskip=8pt},subsection={format=\fontsize{12}{16}\selectfont\bfseries,beforeskip=12pt,afterskip=6pt},subsubsection={format=\normalsize\bfseries,beforeskip=9pt,afterskip=4pt}}
\pagestyle{fancy}\fancyhf{}
\renewcommand{\headrulewidth}{.3pt}
\renewcommand{\footrulewidth}{0pt}
\fancyhead[L]{\parbox[b]{\linewidth}{\fontsize{8}{10}\selectfont\color{Muted}房贷还款渠道：双周扣款与本金入账·第三方付款服务登记及授权条款}}
\fancyfoot[C]{\small\color{Muted}\thepage}
\hypersetup{pdftitle={第三方付款服务登记及授权条款},pdfsubject={仅供模拟使用·房贷还款渠道：双周扣款与本金入账·12-01},pdfauthor={金融决策模拟材料}}
\begin{document}
{\noindent\fontsize{20}{25}\selectfont\bfseries 第三方付款服务登记及授权条款\par}
\vspace{7pt}{\noindent\small\color{Muted}仅供模拟使用·房贷还款渠道：双周扣款与本金入账·12-01\par}\vspace{10pt}
\section*{1. 扣款指示}
服务按登记的付款日历，每隔两周从指定存款账户扣取正常月供的一半。全年日历列示26次半额扣款。每次另收交易费2.50美元。扣款银行、账户号码及付款授权由登记\allowbreak\mbox{资料确定。}


\section*{2. 服务登记费用}
加入服务的登记费用为295美元。登记费用从计划产生的额外半额付款中收取；尚未收足前，额外付款先用于登记费用，余额才按安排\allowbreak\mbox{转付本金。}


\section*{3. 资金保管及月供转付}
从存款账户扣款，并不同时将该笔款项应用于房贷本金。服务暂存所收半额款项，按正常月供日期转付完整月供。贷款机构接收并按贷款账户规则\allowbreak\mbox{应用付款。}


\section*{4. 三次扣款月份}
某月发生三次半额扣款时，超出正常月供所需金额的半额款项为额外付款。按费用收取顺序，先支付尚未收完的登记费用，再把剩余款项用于下一次正常付款安排中的\allowbreak\mbox{本金付款。}


\section*{5. 贷款账户与服务账户}
正常月供义务仍由原贷款合同确定。服务负责收集和转付资金，贷款机构负责贷款账目。所收交易费用与登记费用不作为贷款本金付款，也不改变合同利率或\allowbreak\mbox{贷款期限。}


\section*{6. 扣款授权}
登记申请须载明每次扣款金额、服务费、扣款账户及扣款周期。实际月供到期日及额外本金入账日期以贷款服务机构\allowbreak\mbox{记录为准。}

\end{document}
